\documentclass[12pt]{article}
\usepackage[T1]{fontenc}
\usepackage[utf8]{inputenc}
\usepackage{lmodern}
\usepackage{textcomp}
\usepackage{enumitem}
\usepackage{geometry}
\geometry{margin=1in}
\begin{document}
\begin{center}
Answer Key - Mathematics - Graduate Level Assessment 11 \
\small (Answers are ordered exactly as in the question paper.)
\end{center}

\section*{Multiple Choice}
\begin{enumerate}[label=\arabic*.]
\item \textbf{Answer:} A
\\ \textit{Options:} A) The set of all finite subsets of R; B) R; C) The set of all infinite subsets of R; D) The set of all functions from Z to Z; E) The set of all functions from \{0, 1\} to R
\\ \textit{Note:} The set of all finite subsets of R has the greatest cardinality because it can be shown to have the same cardinality as the power set of R, which is uncountably infinite, thus exceeding the cardinality of any other set of finite subsets.
\item \textbf{Answer:} A
\\ \textit{Options:} A) 2/50; B) 1/25; C) 1/30; D) 3/75; E) 4/100
\\ \textit{Note:} The correct answer of 2/50 arises from applying the hypergeometric distribution to calculate the probability of selecting exactly 2 damaged suitcases out of 3 chosen from a total of 25 suitcases, 5 of which are damaged.
\item \textbf{Answer:} A
\\ \textit{Options:} A) \$1.50; B) \$0.50; C) \$0.20; D) \$0.60; E) \$0.10
\\ \textit{Note:} The correct answer is $0.40, as the unit cost is calculated by dividing the total price of $2.00 by the number of lemons, which is 5, resulting in $2.00 ÷ 5 = $0.40 per lemon.
\item \textbf{Answer:} A
\\ \textit{Options:} A) 5; B) 25; C) 7; D) 6; E) 10
\\ \textit{Note:} The correct answer is 5 because dividing 35 by 5 yields 7, which satisfies the equation 35 / ? = 7.
\item \textbf{Answer:} A
\\ \textit{Options:} A) 24; B) 28; C) 20; D) 16; E) 8
\\ \textit{Note:} In the ring Z\_24, the product (12)(16) equals 192, and when reduced modulo 24, 192 is congruent to 0, which is represented as 24 in the context of the ring.
\end{enumerate}

\section*{True / False}
\begin{enumerate}[label=\arabic*.]
\setcounter{enumi}{5}
\item \textbf{Answer:} True
\\ \textit{Options:} A) True; B) False
\\ \textit{Note:} The statement is true because if seven pens cost $9.24, then the cost per pen is $1.32, making the total cost for three pens \$3.96, which equals the cost of eleven pencils, resulting in each pencil costing 66 cents.
\item \textbf{Answer:} False
\\ \textit{Options:} A) True; B) False
\\ \textit{Note:} The statement is false because the values x = 1, y = 3, and z = -1 do not satisfy all equations in the given system, indicating that they do not represent a valid solution.
\item \textbf{Answer:} True
\\ \textit{Options:} A) True; B) False
\\ \textit{Note:} The statement is true because the number of ways to divide a set of 5 elements into 2 non-empty ordered subsets is calculated as $2^5$ - 2, which equals 30, and then multiplied by 2! (to account for the order of the subsets), resulting in 240 ways.
\item \textbf{Answer:} False
\\ \textit{Options:} A) True; B) False
\\ \textit{Note:} The expression simplifies to \(1990 \times (1991 - 1989) = 1990 \times 2 = 3980\), not 3970, making the statement false.
\item \textbf{Answer:} True
\\ \textit{Options:} A) True; B) False
\\ \textit{Note:} The statement is true because integrating the differential equation dy/dt = √t from t=1 to t=4 yields y(4) = 1 + (2/3)($t^(3$/2)) evaluated from 1 to 4, which approximates to 7.789.
\end{enumerate}

\section*{Long Form Response}
\begin{enumerate}[label=\arabic*.]
\setcounter{enumi}{10}
\item \textbf{Answer:} The division of 2,314 by 4 yields a quotient of 578 and a remainder of 2, denoted as 578 r 2. This result indicates that 2,314 is not evenly divisible by 4, reflecting the concept of modular arithmetic where the remainder provides insight into the relationship between integers. In a broader mathematical context, the quotient represents the largest multiple of 4 that fits within 2,314, while the remainder signifies the excess that cannot be accounted for by these multiples. Such division operations are fundamental in number theory, influencing various applications, including algorithm design and cryptography, where understanding divisibility and remainders plays a crucial role.
\\ \textit{Note:} correct because it accurately computes the quotient and remainder of the division, illustrating the principles of divisibility and modular arithmetic, where the quotient reflects the largest multiple of the divisor within the dividend, and the remainder indicates the surplus that remains after accounting for those multiples.
\item \textbf{Answer:} To find the smallest positive integer that meets the given modular conditions, we can use the method of successive substitutions or the Chinese Remainder Theorem. We need to solve the system of congruences: \( x \equiv 5 \ (\text{mod} \ 8) \), \( x \equiv 1 \ (\text{mod} \ 3) \), and \( x \equiv 7 \ (\text{mod} \ 11) \). Starting with the first two congruences, we can express \( x \) in terms of \( k \) as \( x = 8 k + 5 \), and then substitute into the second congruence to find \( k \). Iterating through values leads to the solution \( x = 103 \) which also satisfies the third condition. Thus, the smallest positive integer that satisfies all three conditions is 103.
\\ \textit{Note:} correct because it systematically applies the method of successive substitutions to solve the system of congruences, beginning with the first two conditions to express \( x \) in terms of one variable, and then substituting this expression into the remaining congruence to find the specific integer that satisfies all conditions.
\end{enumerate}

\vfill
\noindent\textit{End of Answer Key}
\end{document}